% ============================================================
%  II. DCPAM Model Overview
% ============================================================
\section{DCPAM Model Overview}

The DCPAM model takes as input the images from two cameras and
the length of a measurement tool bar, and outputs the
perpendicular distance from the target point to the laser
reference axis:
%
\begin{equation}
  \boxed{h_t = \mathrm{DCPAM}(I_{C_1},\, I_{C_2},\, l_T)}
  \label{eq:dcpam-overview}
\end{equation}
%
where $h_t$ is the target-point-to-axis distance,
$I_{C_1}$ and $I_{C_2}$ are the images captured by cameras
$C_1$ and $C_2$ respectively,
and $l_T$ is the length of the tool bar (side gauge).

Two parameter-tuning strategies are considered:

\textbf{DCPAM-CV.}
An experience-driven workflow in which the operator
iteratively adjusts CV algorithm parameters
(thresholds, filter sizes, fitting methods),
evaluates on a test set, and manually refines the
pipeline until the target accuracy is achieved.

\textbf{DCPAM-CNN.}
A deep-learning approach that trains a convolutional
neural network to directly regress the laser spot
centre from raw images, replacing hand-tuned CV
processing.

% TODO: 补充 DCPAM 整体流程图（建议画一个 pipeline figure）
